\newcommand*{\slcCoverSynopsis}{%
  हे अधितर्कशास्त्राचे परिचयात्मक पाठ्यपुस्तक आहे. यात अनौपचारिक
  संचसिद्धान्त, प्रथम-क्रम तर्कशास्त्र, क्रमवर्ती कलन, नैसर्गिक निगमन,
  संपूर्णता, संहतता व लोव्हेनहाइम--स्कोलेम प्रमेये, ट्यूरिंग यंत्रे,
  तसेच थांबण्याची समस्या व प्रथम-क्रम तर्कशास्त्र यांची अनिर्णेयता
  यांचा अभ्यास केला आहे. \textit{\href{https://forallx.openlogicproject.org/}{forall x}}
  सारख्या ग्रंथातून औपचारिक तर्कशास्त्राची प्राथमिक ओळख झालेल्या
  पदवीच्या विद्यार्थ्यांसाठी हे पुस्तक आहे.}

\renewcommand*{\olremixedby}{%
  \href{\OLPauthorlink}{\OLPauthor}\\
  \href{https://openlogicproject.org/}{Open Logic Project} च्या ग्रंथावर आधारित\\[.6em]
  मराठी यंत्रानुवाद\\[.3em]
  {\small OpenAI Codex (GPT-6.1 Sol, Ultra)}}

\renewcommand*{\olacknowledgements}{%
  University of Calgary येथील
  \href{https://taylorinstitute.ucalgary.ca/}{Taylor Institute of Teaching and Learning}
  आणि Alberta सरकारच्या अर्थसाहाय्याने चालणाऱ्या Alberta Open Educational
  Resources (ABOER) Initiative यांच्या उदार पाठबळाबद्दल
  Open Logic Project त्यांचे आभार मानतो.}

\renewcommand*{\olillustrations}{%
  मुखपृष्ठ आणि भागांच्या प्रारंभीची रेखाचित्रे
  \href{https://github.com/OpenLogicProject/portraits}{Matthew Leadbeater, Open Logic Portraits}
  यांची आहेत (© २०१६). ती
  \href{https://creativecommons.org/licenses/by-nc/4.0/}{CC BY-NC 4.0}
  परवान्याखाली वापरली आहेत. छायाचित्रांचे श्रेय आणि त्यांच्या वापराच्या
  अटी पुस्तकाच्या शेवटी दिल्या आहेत.}

\renewcommand*{\olcolophon}{%
  XeLaTeX मध्ये OL Marathi Serif (Noto Serif Devanagari वर आधारित)
  आणि Open Logic च्या गणितीय अक्षररूपांचा वापर करून अक्षरजुळणी केली आहे.}

\renewcommand*{\olbookrevision}{%
  Richard Zach यांच्या मूळ ग्रंथाची शरद ऋतू २०२५ (F25) आवृत्ती
  (\href{https://github.com/rzach/sets-logic-computation/tree/6178f32f61645f3b62ef621b4c5a88f8e67778ad}{6178f32})
  आणि Open Logic Project चा ग्रंथ
  (\href{https://github.com/OpenLogicProject/OpenLogic/tree/1e960beff9ed7835bf3e3f1335e21af3439cd107}{1e960bef})
  हे स्रोत आहेत. मराठी मजकुरासाठी
  \href{https://github.com/KokunoYumeto/OpenLogic-mr-Deva-IN/releases/tag/complete-v1.2}{\texttt{complete-v1.2}}
  मधील Open Logic चा यंत्रानुवाद वापरला आहे
  (OpenAI Codex, GPT-5.6 Sol व GPT-6 Sol, दुरुस्त्यांचे परीक्षण GPT-6.1 Sol).
  या पाठ्यपुस्तकासाठीचा मराठी यंत्रानुवाद आणि दुरुस्त्या
  OpenAI Codex, GPT-6.1 Sol यांनी केल्या आहेत. सर्व मॉडेलांसाठी
  तर्कप्रयत्नाची पातळी Ultra होती.}

\renewcommand*{\ollicensetext}[3]{%
  #2 यांच्या \textit{#1} या ग्रंथाचा मजकूर #3 परवान्याखाली उपलब्ध आहे.\space}

\renewcommand*{\oluselicensetext}{%
  हा ग्रंथ \href{https://openlogicproject.org/}{Open Logic Project} च्या
  \href{https://github.com/OpenLogicProject/OpenLogic/tree/1e960beff9ed7835bf3e3f1335e21af3439cd107}{\textit{The Open Logic Text}} वर आधारित आहे.
  त्या ग्रंथाचा मजकूरही
  \href{https://creativecommons.org/licenses/by/4.0/}{CC BY 4.0}
  परवान्याखाली वापरला आहे.}
